Among outlets recorded as attending the same press conference, those that ask
a question have higher pre-lottery ratios of questions to recorded attendance.
The estimated differences range from 0.152 to 0.185 questions per attendance
record across the four historical windows; the individual 95\% confidence
intervals exclude zero in every window. These associations are consistent with
persistent differences in engagement with the question-and-answer process.
One possible explanation is that some outlets attend primarily to seek an
opportunity to question the president, whereas others attend primarily to
cover the proceedings. The data do not directly measure these intentions,
however, and prior participation may also reflect differences in access,
resources, or editorial practices.

The joint tests provide a complementary diagnostic. Conditional on conference-day
fixed effects and indicators for unavailable histories, the twelve historical
measures jointly predict questioning for the preceding 18- and 24-month windows,
with Holm-adjusted $p$-values of 0.014 and 0.002. The tests for the preceding
12 months and the full pre-lottery period do not reject at the 5\% level.
These omnibus results cannot be attributed specifically to participation
history, because the tested set also includes contracts and YouTube measures.
They are descriptive diagnostics: YouTube engagement is measured from saved
snapshots of videos published before the lottery, not necessarily from counts
collected before the lottery itself.

Taken together, these patterns motivate an outcome analysis that accounts for
pre-lottery participation histories rather than treating access to questions
as unconditionally random. Such adjustment can allow outlets with different
prior engagement to have different untreated outcome paths. It does not, by
itself, establish that question access is conditionally random or remove
selection on unobserved characteristics. A conditional causal interpretation
therefore requires additional identifying assumptions and supporting diagnostics.

% Implementation reminder (not typeset): frozen outlet histories are absorbed
% by outlet fixed effects in the existing FEct specification. Simply adding
% their levels does not adjust the model further. A prespecified interaction
% with calendar time, or another explicit adjustment design, must be agreed
% before re-estimation. Do not use snapshot engagement as a verified baseline.
% Individual tests use available-history samples; the joint tests instead
% retain all 81 outlet units using missingness indicators. Participation-ratio
% confidence intervals are individual, not simultaneous/multiplicity-adjusted.
